\documentclass[11pt,letterpaper,spanish]{article}

\usepackage{caption}

\usepackage{iftex}
\ifPDFTeX
  % --- Setup for pdfLaTeX ---
  \usepackage[utf8]{inputenc}
  \usepackage[T1]{fontenc}
\else
  % --- Setup for XeLaTeX & LuaLaTeX ---
  \usepackage{fontspec}
  \setmainfont{Latin Modern Roman}
\fi


\usepackage[spanish]{babel}

\usepackage{url}
\usepackage{hyperref}
\usepackage{geometry}
\usepackage{svg}
\usepackage{fancyhdr}
\usepackage{xcolor}
\usepackage{multicol}
\usepackage{enumitem}

%% Paquetes caseros
\usepackage{customcode}

% ==========================================
% MARCADO DE PROCEDENCIA DE LAS JUSTIFICACIONES
% (permite distinguir lo tomado de Explicación.pdf
%  de las propuestas de mejora para poder depurarlas)
% ==========================================
\definecolor{propcolor}{RGB}{0,90,160}

\tcbuselibrary{breakable}

\newcommand{\marcaOriginal}{}
\newcommand{\marcaPropuesta}{\textnormal{\small\textcolor{propcolor}{(propuesta)}}}

\newtcolorbox{cajapropuesta}{
	colback=propcolor!4,
	colframe=propcolor,
	boxrule=0.5pt,
	arc=2pt,
	left=5pt, right=5pt, top=5pt, bottom=5pt,
	breakable
}



% Required by biblatex-apa for context-sensitive quotes
\usepackage{csquotes}

% Load biblatex with official APA 7th edition style
\usepackage[style=apa, backend=biber]{biblatex}
\addbibresource{referencias.bib}


\geometry{margin=2cm, top=2.5cm, bottom=2.5cm}

\definecolor{headergray}{gray}{0.3}
\definecolor{rulecolor}{gray}{0.6}

\pagestyle{fancy}
\fancyhf{}
\renewcommand{\headrulewidth}{0.5pt}
\renewcommand{\footrulewidth}{0.5pt}
\fancyhead[L]{\small\textcolor{headergray}{UNIVERSIDAD NACIONAL AUTÓNOMA DE MÉXICO}}
\fancyhead[R]{\small\textcolor{headergray}{\documentoTitulo}}
\fancyfoot[L]{\small\textcolor{headergray}{FACULTAD DE CIENCIAS, 2027-I}}
\fancyfoot[C]{\small\textcolor{headergray}{\thepage}}
\fancyfoot[R]{\small\textcolor{headergray}{FUNDAMENTOS DE BASES DE DATOS}}

% ==========================================
% CONFIGURACIÓN DE VARIABLES (EDITAR AQUÍ)
% ==========================================
\newcommand{\documentoTitulo}{Práctica 05}
\newcommand{\documentoSubtitulo}{Lenguaje para Definición de Datos DDL}

\newcommand{\profesorNombre}{Gerardo Avilés Rosas}

\newcommand{\ayudanteTeoriaUno}{Oscar Daniel Flores Linares}
\newcommand{\ayudanteTeoriaDos}{Claudia Itzel Gutiérrez Sánchez}

\newcommand{\ayudanteLabUno}{Ricardo Badillo Macías}
\newcommand{\ayudanteLabDos}{Rocío Aylin Huerta González}

\newcommand{\equipoNombre}{Dobles Comillas}

\newcommand{\integranteUno}{Contreras Morales Román Ismael}
\newcommand{\integranteDos}{Cornejo de la mora Iñaki}
\newcommand{\integranteTres}{Hipólito Felipe Christopher Guillermo}
\newcommand{\integranteCuatro}{Ledesma Cuevas Eduardo}

% ==========================================

\begin{document}

\thispagestyle{empty}

\begin{center}
	\vspace*{1cm}

	\begin{tabular}{cc}
		\includegraphics[width=3cm]{escudos/unam_seal.pdf} & \includegraphics[width=3cm]{escudos/fc_seal.pdf}
	\end{tabular}

	\vspace{0.8cm}
	{\LARGE\bfseries UNIVERSIDAD NACIONAL AUTÓNOMA DE MÉXICO}\\[0.3cm]
	{\large FACULTAD DE CIENCIAS, 2027-I}\\[0.3cm]
	{\large\bfseries FUNDAMENTOS DE BASES DE DATOS}

	\vspace{1.2cm}
	{\huge\bfseries \documentoTitulo:}\\[0.3cm]
	{\Large \documentoSubtitulo}

	\vspace{1.5cm}
\end{center}

\begin{multicols}{2}
	\raggedcolumns

	\textbf{PROFESOR:}\\
	\profesorNombre\\[1.2cm]

	\textbf{AYUDANTES DE TEORÍA:}\\
	\ayudanteTeoriaUno\\
	\ayudanteTeoriaDos\\[1.2cm]

	\textbf{AYUDANTES DE LABORATORIO:}\\
	\ayudanteLabUno\\
	\ayudanteLabDos\\[1.2cm]

	\columnbreak

	\textbf{EQUIPO:}\\
	\equipoNombre\\[1.2cm]

	\textbf{INTEGRANTES DEL EQUIPO:}\\
	\integranteUno\\
	\integranteDos\\
	\integranteTres\\
	\integranteCuatro

\end{multicols}

\vspace{1cm}

\section*{Restricciones de Dominio}

\subsection*{Cliente}
\begin{itemize}
    \item \textbf{Columnas: Nombre, Paterno.} Pertenecen al dominio de cadena de caracteres de longitud variable \texttt{varchar(100)}. Cuentan con una restricción explícita de no nulidad mediante \texttt{NOT NULL}, lo cual asegura que el sistema rechace cualquier inserción que intente dejar el nombre o apellido paterno en blanco.
    \item \textbf{Columna: Materno.} Pertenece al dominio de cadena de caracteres \texttt{varchar(100)}. Estructuralmente carece de la restricción \texttt{NOT NULL} (acepta valores nulos) pues no todos los clientes poseen legalmente un segundo apellido.
    \item \textbf{Columna: FechaNacimiento.} Pertenece al dominio temporal \texttt{date}. Posee la restricción \texttt{NOT NULL}, forzando al sistema a capturar esta información obligatoriamente.
    \item \textbf{Columna: Sexo.} Pertenece al dominio de cadena de caracteres \texttt{varchar(50)}. No posee restricciones adicionales de nulidad por preferencia del cliente.
\end{itemize}

\subsection*{Telefono\_Cliente}
\begin{itemize}
    \item \textbf{Columna: idCliente.} Pertenece al dominio de los enteros (\texttt{int}). Al provenir de la entidad fuerte, actúa como llave foránea (FK) referenciando a la relación Cliente, estando sujeta a una estricta restricción de integridad referencial. En conjunto con el número telefónico, conforma la llave primaria de la tabla.
    \item \textbf{Columna: telefono.} Pertenece al dominio de cadena de caracteres de longitud variable \texttt{varchar(15)}. Cuenta con restricción de integridad de dominio explícita (\texttt{NOT NULL}), lo cual impide que el sistema registre cadenas vacías para los medios de contacto.
\end{itemize}

\subsection*{Correo\_Cliente}
\begin{itemize}
    \item \textbf{Columna: idCliente.} Pertenece al dominio de los enteros (\texttt{int}). Actúa como llave foránea (FK) referenciando a la relación Cliente con integridad referencial estricta. En combinación con el correo electrónico, conforma la llave primaria de esta relación.
    \item \textbf{Columna: correo.} Pertenece al dominio de cadena de caracteres \texttt{varchar(100)}. Posee una restricción explícita de integridad de dominio (\texttt{NOT NULL}), garantizando que no se puedan registrar cadenas en blanco.
\end{itemize}

\subsection*{Tarjeta}
\begin{itemize}
    \item \textbf{Columna: idCliente.} Pertenece al dominio de los enteros (\texttt{int}). Actúa como llave foránea (FK) hacia la tabla Cliente. Además, compone la primera mitad de la llave primaria de esta tabla.
    \item \textbf{Columna: idTarjeta.} Pertenece al dominio de los enteros (\texttt{int}). Cuenta con restricción de no nulidad (\texttt{NOT NULL}) y funciona como el discriminante que completa la llave primaria compuesta.
    \item \textbf{Columna: Saldo.} Se define como un valor numérico exacto bajo el dominio \texttt{decimal(8,2)}. Posee una restricción de integridad de dominio (\texttt{NOT NULL}), garantizando que ninguna tarjeta se registre sin un balance inicial, incluso si es cero.
    \item \textbf{Columna: rfc\_cajero.} Pertenece al dominio de cadena de caracteres \texttt{varchar(100)}. Funciona como llave foránea (FK) hacia la tabla Cajero, garantizando mediante integridad referencial que toda tarjeta esté asociada al empleado válido que la emitió.
\end{itemize}

\subsection*{Recarga}
\begin{itemize}
    \item \textbf{Columna: idRecarga.} Pertenece al dominio de los enteros (\texttt{int}). Opera con restricción de integridad de entidad (\texttt{NOT NULL}, \texttt{UNIQUE}) al funcionar como la llave primaria de la tabla.
    \item \textbf{Columnas: MontoAbonar, fecha, hora, MetodoPago.} Operan bajo los dominios numérico exacto (\texttt{decimal(8,2)}), temporal de fecha (\texttt{date}), temporal de hora (\texttt{time}) y cadena de caracteres (\texttt{varchar(50)}) respectivamente. Todas cuentan con una restricción explícita de integridad de dominio (\texttt{NOT NULL}), asegurando que no existan transacciones financieras sin los datos obligatorios.
    \item \textbf{Columna: idSucursal.} Pertenece al dominio \texttt{int}. Funciona como llave foránea (FK) aplicando integridad referencial hacia la relación Sucursal.
    \item \textbf{Columnas: idCliente, idTarjeta.} Pertenecen al dominio \texttt{int}. Conforman una llave foránea (FK) compuesta que mantiene la integridad referencial hacia la llave primaria compuesta de la relación Tarjeta.
    \item \textbf{Columna: rfc\_cajero.} Pertenece al dominio \texttt{varchar(100)}. Actúa como llave foránea (FK) con integridad referencial hacia la relación Cajero.
\end{itemize}

\subsection*{Sucursal}
\begin{itemize}
    \item \textbf{Columna: idSucursal.} Pertenece al dominio de los enteros (\texttt{int}). Aplica una restricción de integridad de entidad (\texttt{NOT NULL}, \texttt{UNIQUE}), estableciéndose como la llave primaria de la relación.
    \item \textbf{Columnas: Nombre, Calle, Colonia, Estado, Horario.} Pertenecen al dominio de cadena de caracteres de longitud variable \texttt{varchar(100)}. El conjunto cuenta con restricción explícita de integridad de dominio (\texttt{NOT NULL}), forzando el registro detallado de la ubicación e identidad de la sucursal.
    \item \textbf{Columnas: NumExterior, NumInterior.} Pertenecen al dominio de cadena de caracteres \texttt{varchar(50)} y cuentan con restricción explícita de integridad de dominio (\texttt{NOT NULL}).
    \item \textbf{Columna: Telefono.} Pertenece al dominio de cadena de caracteres \texttt{varchar(15)} y cuenta con restricción explícita de integridad de dominio (\texttt{NOT NULL}).
\end{itemize}

\subsection*{Maquina}
\begin{itemize}
    \item \textbf{Columna: idMaquina.} Pertenece al dominio de los enteros (\texttt{int}). Cuenta con restricción de integridad de entidad (\texttt{NOT NULL}, \texttt{UNIQUE}) al ser la llave primaria de la tabla.
    \item \textbf{Columna: Estado.} Pertenece al dominio de cadena de caracteres \texttt{varchar(50)}. Aplica una restricción de integridad de dominio explícita (\texttt{NOT NULL}) pues una maquina tiene que estar en un estado.
    \item \textbf{Columna: idSucursal.} Pertenece al dominio de los enteros (\texttt{int}). Ejerce como llave foránea (FK) con integridad referencial hacia la relación Sucursal. Cuenta con restricción estricta de no nulidad (\texttt{NOT NULL}) pues una maquina debe estar en una sucursal.
    \item \textbf{Columna: idJuego.} Pertenece al dominio de los enteros (\texttt{int}). Funciona como llave foránea (FK) manteniendo integridad referencial hacia la relación Juego, bajo una restricción estricta de no nulidad (\texttt{NOT NULL}) pues una maquina esta asociada a un juego.
\end{itemize}

\subsection*{Juego}
\begin{itemize}
    \item \textbf{Columna: idJuego.} Pertenece al dominio de los enteros (\texttt{int}). Mantiene una restricción de integridad de entidad (\texttt{NOT NULL}, \texttt{UNIQUE}), consolidándose como la llave primaria de la tabla.
    \item \textbf{Columnas: Nombre, Tipo.} Pertenecen al dominio de cadena de caracteres de longitud variable \texttt{varchar(100)} y \texttt{varchar(50)} respectivamente. Ambas cuentan con una restricción explícita de integridad de dominio (\texttt{NOT NULL}).
    \item \textbf{Columna: Costo.} Pertenece al dominio numérico exacto \texttt{decimal(8,2)}. Cuenta con restricción de integridad de dominio (\texttt{NOT NULL}), asegurando que no existan juegos sin un costo financiero asignado.
\end{itemize}

\subsection*{Partida}
\begin{itemize}
    \item \textbf{Columna: idPartida.} Pertenece al dominio de los enteros (\texttt{int}). Actúa como la llave primaria de la relación, por lo que aplica una restricción de integridad de entidad estricta (\texttt{NOT NULL}, \texttt{UNIQUE}).
    \item \textbf{Columnas: Puntaje, PuntosGanados.} Pertenecen al dominio de los enteros (\texttt{int}). Cuentan con una restricción explícita de integridad de dominio (\texttt{NOT NULL}), asegurando el registro cuantitativo de la interacción del cliente.
    \item \textbf{Columnas: fecha, hora.} Operan bajo los dominios temporales de fecha (\texttt{date}) y hora (\texttt{time}) respectivamente. Ambas están protegidas con restricciones de integridad de dominio (\texttt{NOT NULL}).
    \item \textbf{Columna: idMaquina.} Pertenece al dominio de los enteros (\texttt{int}). Absorbe la llave de la relación Máquina, fungiendo como llave foránea (FK) con integridad referencial.
    \item \textbf{Columnas: idCliente, idTarjeta.} Pertenecen al dominio de los enteros (\texttt{int}). Conforman una llave foránea (FK) compuesta que mantiene la integridad referencial apuntando exactamente a la llave primaria compuesta de la relación Tarjeta.
\end{itemize}

\subsection*{Canje}
\begin{itemize}
    \item \textbf{Columna: idCanje.} Pertenece al dominio de los enteros (\texttt{int}). Está sujeta a la restricción de integridad de entidad (\texttt{NOT NULL}, \texttt{UNIQUE}) por ser la llave primaria de la tabla.
    \item \textbf{Columnas: fecha, hora, PuntosUtilizados, CantidadPremiosObtenidos.} Operan bajo los dominios temporales \texttt{date} y \texttt{time}, y el dominio numérico entero \texttt{int} respectivamente. Todas operan bajo restricciones explícitas de integridad de dominio (\texttt{NOT NULL}), forzando al sistema a registrar los detalles precisos de cada transacción.
    \item \textbf{Columnas: idSucursal, idCliente, idPremio.} Pertenecen al dominio de los enteros (\texttt{int}). Actúan como llaves foráneas (FK) individuales que aplican reglas de integridad referencial hacia las tablas Sucursal, Cliente y Premio.
    \item \textbf{Columna: rfc\_EncargadoDePremios.} Pertenece al dominio de cadena de caracteres \texttt{varchar(100)}. Funge como llave foránea (FK) con integridad referencial hacia la subentidad EncargadoDePremio, vinculando al empleado responsable de la autorización.
\end{itemize}

\subsection*{Premio}
\begin{itemize}
    \item \textbf{Columna: idPremio.} Pertenece al dominio de los enteros (\texttt{int}). Ostenta la restricción de integridad de entidad (\texttt{NOT NULL}, \texttt{UNIQUE}) al ser la llave primaria del catálogo.
    \item \textbf{Columnas: Nombre, RangoEdad, Categoria.} Pertenecen al dominio de cadena de caracteres \texttt{varchar(100)}, \texttt{varchar(50)} y \texttt{varchar(50)} respectivamente. Cuentan con restricciones explícitas de integridad de dominio (\texttt{NOT NULL}), asegurando que todo premio esté correctamente clasificado.
    \item \textbf{Columna: PuntosNecesarios.} Pertenece al dominio de los enteros (\texttt{int}). Protegida con restricción de no nulidad (\texttt{NOT NULL}) para evitar que un premio se registre sin su costo en puntos.
    \item \textbf{Columna: ValorAprox.} Pertenece al dominio numérico exacto \texttt{decimal(8,2)}. Cuenta con restricción de integridad de dominio (\texttt{NOT NULL}).
\end{itemize}


\printbibliography

\end{document}

